% ============================================================================
\section{Related Work}
\label{sec:related}
% ============================================================================

\paragraph{Neural Scene Relighting.}
\todo{}
NeRF-OSR~\cite{rudnev2022nerfOSR} decomposes outdoor NeRF scenes into albedo and per-image SH lighting, enabling appearance transfer between views.
NeRF-W~\cite{martinbrualla2021nerfw} handles appearance variation with per-image latent codes.
More recently, methods based on 3D Gaussian Splatting~\cite{kerbl20233dgs} have achieved real-time rendering while maintaining relighting capability.

\paragraph{Outdoor Illumination Modeling.}
\todo{}Modeling outdoor illumination typically involves representing the sun as a directional source and the sky as a hemisphere of diffuse light~\cite{lalonde2012estimating,hold-geoffroy2019deep}.
The color of sunlight varies with elevation due to atmospheric Rayleigh scattering, producing the characteristic warm tones at sunrise and sunset~\cite{preetham1999practical,nishita1993display}.
Our sun color prior draws on this atmospheric model.

\paragraph{Shadow Computation.}
\todo{}Shadow mapping~\cite{williams1978casting} renders a depth map from the light's viewpoint and compares fragment depths to determine occlusion.
In the Gaussian splatting setting, the unstructured point-based representation requires adapting shadow mapping to work with splatted depth rather than mesh rasterization.

\paragraph{2D Gaussian Splatting.}
\todo{}2DGS~\cite{huang20242dgs} represents scene geometry as oriented planar splats (surfels), providing well-defined surface normals from the splat orientation.
LumiGauss~\cite{joaxkal2024lumigauss} builds on 2DGS, leveraging the surfel normals for hemisphere-based SH evaluation.
We inherit this representation and additionally leverage the normals for Lambert shading (N$\cdot$L).